\begin{textsample}{POS Dim 3 – am – Score 36.00 – am/am\_inf\_39.txt}  \label{ex:f3_pos_034}
18.3 O ambiente educativo e seus diferentes espaços na fase \textbf{creche} e pré-escola Ao pensar no processo de construção de espaços para \textbf{bebês}, \textbf{crianças} bem \textbf{pequenas} e \textbf{criança} \textbf{pequena} da Educação \textbf{Infantil}, precisamos compreender a forma como se dá o processo de desenvolvimento e aprendizagem.
Desse modo, cabe a/ao professora/professor planejar atividades experienciais que promovam desafios, tendo o espaço como um colaborador fundamental nesse processo.
É preciso considerar a organização dos tempos e espaços das unidades de ensino da Educação \textbf{Infantil} dos municípios do Amazonas e compor um cenário educativo integrador entre criança/criança, adulto/criança e criança/comunidade, num contexto em que as relações sociais e a cultura sejam parte do processo educativo.
Os recursos utilizados para o desenvolvimento da prática pedagógica com \textbf{bebês}, \textbf{crianças} bem \textbf{pequenas} e \textbf{crianças} \textbf{pequenas} devem ser adequados e seguros para \textbf{bebês} e \textbf{crianças}, num espaço que provoque o interesse de \textbf{manipular} e conhecer a variedade de materiais e objetos da cultura e que crie condições para o exercício da liberdade, construindo, assim, sua identidade pessoal e coletiva, por meio das diversas vivências.
O espaço precisa ser um ambiente dialógico, que educa, contemplando as diferentes culturas, povos, raças e modos de vida.
Além dos materiais e \textbf{brinquedos} industrializados, é importante que haja \textbf{brinquedos} não estruturados e materiais reciclados seguros.
A organização dos espaços para atendimento dos \textbf{bebês}, \textbf{crianças} bem \textbf{pequenas} e \textbf{crianças} \textbf{pequenas} de até 5 anos e 11 \textbf{meses} de idade deve estar em consonância com o que preconizam os Parâmetros Básicos de Infraestrutura das \textbf{Instituições} de Educação \textbf{Infantil} ( BRASIL, 2006 ) quando sugerem ambientes destinados ao \textbf{cuidar}, educar e \textbf{brincar}, buscando integrar o espaço físico da unidade de Educação \textbf{Infantil} às atividades experienciais ali desenvolvidas.
Esses espaços devem ser preparados, considerando os \textbf{bebês} e as \textbf{crianças} em uma relação intersubjetiva, \textbf{atentando} ao ritmo próprio de cada um, que necessita dentre outras coisas, realizar a prática da descoberta por meio do engatinhar, rolar, \textbf{apoiar} -se e assim, preparar -se para os primeiros passos.
Nesse momento, a/o professora/professor precisa estimular o \textbf{bebê} e a \textbf{criança} bem \textbf{pequena} a explorar diversos tipos de materiais, observar, \textbf{brincar}, tocar, alimentar -se, fazer sua \textbf{higiene} pessoal, descansar e dormir, realizando, assim, suas necessidades básicas e sociais. \textbf{Bebês} e \textbf{crianças} devem estar em espaços ricamente planejados e intencionalmente organizados para favorecer a exploração físico-motor e a capacidade de estabelecer relações com os espaços que a rodeia, bem como agir sobre diferentes ambientes construídos a partir de sua atuação como protagonista do seu fazer.
A ambientação da sala deve ser composta de atividades planejadas com e para os \textbf{bebês} e as \textbf{crianças}, de forma intencional.
As múltiplas atividades desenvolvidas e criadas no ambiente educacional não \textbf{doméstico} inserem os \textbf{bebês} e as \textbf{crianças} em um contexto instigante, que a faça se sentir afetada pelos objetos da cultura e as relações com os outros, enriquecendo e ampliando as possibilidades de entendimento de si e do mundo que as cercam.
O espaço externo da \textbf{instituição} deve ser explorado tanto quanto o interno, o que contribui para que \textbf{bebês} e \textbf{crianças} construam as primeiras noções de \textbf{cuidado} com o meio ambiente, interagindo com \textbf{adultos} e \textbf{crianças} de diferentes idades.
As \textbf{brincadeiras} realizadas neste meio visam a exploração desse espaço, garantindo o contato dos \textbf{bebês} e \textbf{crianças} com a natureza, como também com os habitantes desse espaço, como formigas, minhocas, passarinhos, tartarugas, jabutis, galinhas, peixes, dentre outros.
Deve -se ainda considerar, quando houver, o espaço externo da \textbf{instituição} como uma área de possibilidades de aprendizagem e desenvolvimento, proporcionando às \textbf{crianças} vivências em um espaço livre, experiências de \textbf{movimentação} ampla, de convivência com \textbf{brincadeiras} no mundo físico, experienciando a terra, a água e a relação com seus \textbf{pares}.
As interações desenvolvidas nestes espaços são ricas e desafiadoras, especialmente quando há plantas, terra, \textbf{brinquedos} para \textbf{brincadeiras} coletivas, como balanço, cordas, bola etc., e espaço para produzir arte que com materiais como argila, tintas corporais, \textbf{pintura} na \textbf{parede}, entre outras.
A construção desse ambiente externo se dá também com a participação da família e comunidade.
Elas podem trazer para a escola plantas, sementes, participar da produção de cartazes, do plantio de mudas e assim, se integrar ao ambiente institucional, ampliando sua visão de mundo associada a dos \textbf{bebês} e \textbf{crianças} e, sobretudo, proporcionar uma relação viva e de conhecimento entre família, comunidade e \textbf{instituição}, construindo uma educação de qualidade e de referência.
Os Parâmetros Básicos de Infraestrutura das \textbf{Instituições} de Educação \textbf{Infantil} ( BRASIL, 2006b ) orientam que esse espaço deve contemplar chuveiros, torneiras acessíveis ao tamanho das diferentes faixas \textbf{etárias}, quadros de azulejos laváveis para atividades com tinta à base d’água, \textbf{brinquedos} de \textbf{parque}, como escorregador, trepa-trepa, balanços, túneis etc., contendo diversos tipos de pisos, como : grama, terra e área cimentada.
Esta área deve \textbf{contar}, ainda, com espaços ensolarados e sombreados, com áreas verdes que podem \textbf{contar} com pomar, horta e jardim, proporcionando ao \textbf{bebê} e à \textbf{criança}, aprendizagem e desenvolvimento integrador.
Deve -se considerar os diferentes espaços como um ambiente integrador e facilitador no processo de inclusão que permite à \textbf{criança} interagir com o outro e ter acesso aos bens culturais, o que pode acontecer pelas zonas circunscritas.
As zonas circunscritas⁵ e suas organizações nas salas de referências possibilitam aos \textbf{bebês} e \textbf{criança}, diferentes formas de atuação e de vivência com o faz de conta e por meio das interações e \textbf{brincadeiras}, propicia a si e ao outro, novas formas experienciais de compreender os diferentes modos de vida, cultura, costumes e saberes diferentes do seu, contribuindo ainda com um olhar investigativo sobre as coisas e os objetos ali inseridos.
As zonas circunscritas podem ser construídas de diferentes formas e maneiras, isto vai depender do espírito investigador da/do professora/professor em perceber nas \textbf{crianças} seus anseios e \textbf{desejos} na descoberta do novo e assim construir a partir do olhar da turma, um planejamento que contemple essa construção.
Por isso, as zonas circunscritas e os cantos não devem ser pensados como algo fixo, mas pensados e repensados, construídos e reconstruídos conforme o interesse e necessidade dos \textbf{bebês} e \textbf{crianças}.
A seguir, algumas \textbf{sugestões} que podem compor as Zonas Circunscritas e/ou cantos.
Lembrando que a ampliação e a criação de outros espaços dependem daquilo que os \textbf{bebês} e \textbf{crianças} sinalizam como Zonas Circunscritas. | Zonas Circunscritas/Cantos | Materiais | |---|---| | Zonas/Cantos da casa de bonecas | Fogão a lenha e convencional, geladeira, TV, cama, mesa, outros móveis, objetos para utensílios de cozinha, quarto, banheiro, sala ( poderão ser confeccionados com material de sucata e com a participação das \textbf{crianças} ) etc. | | Zonas/Cantos da fantasia | Pedaços de \textbf{tecidos} variados, fantasias, chapéus, sapatos, roupas, \textbf{espelho}, maquiagens, óculos, bolsas, cintos, gravatas, bijuterias, artefatos de cabelo ( prendedores, tiaras, grampos, fitas ) etc. | | Zonas/Cantos da \textbf{biblioteca} | Almofadas, tapete, estantes, painel de informações, \textbf{livros}, revistas e jornais. | | Zonas/Cantos da garagem | Tacos de madeira para construção, \textbf{carros}, caminhões, trilhas, placas de sinais de trânsito, cordas, ferramentas e equipamentos para \textbf{manutenção} de \textbf{carros}. | | Zonas/Cantos de jogos e \textbf{brinquedos} | Jogos de encaixe, de armar, de quebra-cabeça, sucatas variadas e organizadas separadamente, peças de madeira. | | Zonas/Cantos de música | Instrumentos musicais comprados ou confeccionados pela/pelo professora/professor e/ou pelas \textbf{crianças}, rádio, aparelho de \textbf{som} com CDs e DVDs variados. | | Zonas/Cantos do supermercado | \textbf{Embalagens} vazias de diferentes produtos ( caixas, garrafas e latas ), frutas e legumes de \textbf{plástico}, sacos para empacotar, caixa registradora, dinheiro de papel, cartazes com nome dos produtos, prateleiras. | | Zonas/Cantos do conto e reconto | \textbf{Livros} de todo o tipo e tamanho, preferencialmente com formas grandes e com poucas palavras.
Conter almofadas, tapetes para uma boa interação e acolhimento, fantoches, dedoches, teatro de sombras. | ⁵ As zonas circunscritas ( ou os cantos de referência ou cantinhos de faz de conta, onde são disponibilizados diferentes tipos de materiais, mas que não direcione a ação da \textbf{criança} ) devem instigar os \textbf{bebês} e \textbf{crianças} à exploração dos espaços pedagógicos.
Segundo Campos de Carvalho e Padovani ( 2000 ), que pesquisam o ato de ressignificar a organização dos espaços nas salas de referências, inserindo a criação de zonas circunscritas, é necessário transformar o ambiente com áreas que propiciem a interação e a \textbf{brincadeira} dos \textbf{bebês} e \textbf{crianças}.

% matched lemmas: adulto, apoiar, atentar, bebê, biblioteca, brincadeira, brincar, brinquedo, carro, contar, creche, criança, cuidado, cuidar, desejo, doméstico, embalagem, espelho, etário, higiene, infantil, instituição, livro, manipular, manutenção, movimentação, mês, par, parede, parque, pequeno, pintura, plástico, som, sugestão, tecido
\end{textsample}
